\documentclass{bareminimum}

\topic{Set Theory}
\category{Mathematics}
\tags{Foundations of Mathematics, Set Theory}

\begin{document}

\maketitle
	
\section{Introduction}

Set theory is important for those working on Assumptions of Physics for at least three reasons. First, it is a foundational framework in mathematics. Second, it showcases a successful attempt at concept generalization. Third, it shows how a formal system is bootstrapped.

There are roughly two branches of set theory: naive and axiomatic. Naive set theory is built on intuitive concepts, and as such is not fully formalized and is open to potential paradoxes. Axiomatic set theory provides a fully formalized axiomatic system that aims to close those problems. We will start with naive set theory, which is the best setting for defining all the basic notions and getting a sense of how the framework works. Then we will turn our attention to axiomatic set theory to get a sense of what problems it solves and how, and whether those solutions mesh with foundational goals in physics.

For more details on set theory, see for example \cite{pinter2014book}.

\end{document}